% This is a template for your written document.
%
% To compile using latexmk on the command line, run the following: 
% latexmk -pdf main.tex

\documentclass[12pt]{article}
\usepackage{setspace}
\singlespace
\usepackage[left=1in,right=1in,top=1in,bottom=1in]{geometry}
\usepackage{graphicx}

\title{\textbf{Holodeck}}
\author{Leslie Osei-Anane}

\begin{document}

\maketitle

\section{Introduction}
Science fiction has always been reliable at predicting what we want before we know how to build it. The communicator from Star Trek became the first commercial flip phone. The tablet computer from 2001: A Space Odyssey became the iPad. These weren’t coincidences. They were the same human desire, expressed first in story, then in silicon.

Of all the technologies science fiction imagined, the Holodeck stuck around the longest without a real answer. You walked in, described a world, and it built one around you. No menus, no asset browsers, no friction between imagination and experience. It felt further away than flying cars because the problem it was solving was harder to name. The hard part was never the hardware. Understanding what a person meant, building something coherent from it, and delivering it before the moment of anticipation turned into waiting was the actual problem.
That sequence of steps is deceptively demanding. A user speaks a sentence. The system interprets intent, selects a layout, furnishes it with appropriate objects, themes it correctly, and verifies that every surface is safe to walk on. In VR specifically, all of that has to work before the user takes a single step. The system earns the experience before it offers it.

Recent work has made real progress on pieces of this. ProcTHOR demonstrates that large-scale valid indoor environments can be generated procedurally with measurable generalization benefits for embodied AI~\cite{Deitke_2022_NeurIPS}. Text2Room shows that a text prompt can seed a spatially consistent 3D mesh by lifting 2D generative outputs into a 3D representation~\cite{hoellein2023text2room}. Work directly named Holodeck has explored language-guided 3D environment generation by selecting and placing assets to satisfy language-specified constraints~\cite{Yang_2024_CVPR}. Each of these advances a different part of the problem. None of them put a user inside a headset and asked the system to perform in real time.
\section{Background}
\subsection{Large Language Models and Structured Output}
Large language models (LLMs) are neural networks trained on large corpora of text to model and generate natural language. When prompted with appropriate constraints, these models produce outputs that conform to structured formats, including JSON schemas, enabling their integration as planning components in software pipelines. Schema-locking, which constrains model output to a predefined JSON schema, reduces the likelihood of hallucinated or invalid field values and allows downstream components to parse and validate model responses deterministically. This property is central to the system described in this paper, which uses an LLM to convert natural language prompts into a validated WorldSpec representation rather than generating 3D content directly. A deterministic fallback planner handles cases where the LLM produces invalid or out-of-registry outputs, ensuring the pipeline succeeds even under model failure.

\subsection{Virtual Reality and OpenXR}
Virtual reality systems present stereoscopic rendered environments through a head-mounted display, tracking head and controller motion to update the view in real time. The OpenXR standard defines a cross-platform API for extended reality applications, allowing a single codebase to target multiple headset platforms. The Meta Quest line of headsets runs OpenXR natively and imposes strict performance constraints: frame time budgets, draw call limits, and texture memory ceilings must be respected to maintain visual stability and user comfort. Teleport locomotion, a standard interaction pattern in VR, allows a user to select a destination on a valid floor surface and reposition instantly, sidestepping continuous motion that can cause discomfort. These constraints shaped the design of this system at every level, from asset budget enforcement to the decision to prioritize greybox compilation before visual quality.

\subsection{Asset-Based Scene Composition}
Asset-based scene composition constructs environments by selecting and placing pre-authored 3D objects from a controlled library, rather than generating geometry procedurally from noise functions or lifting it from 2D outputs. This approach trades generative range for reliability: the resulting environments contain assets with known geometry, collision data, and performance characteristics that a runtime can depend on. Procedural placement systems arrange these assets according to spatial constraints and layout rules while preserving physical validity. This system adopts asset-based composition as its primary generation strategy, supplemented by procedurally generated greybox geometry for structural layout. Assets enter the planner-safe inventory only after passing a manual review pipeline that validates semantic quality and assigns exclusion flags where appropriate.

\subsection{Compiler Pipelines and Intermediate Representations}
A compiler pipeline transforms a high-level description into an executable form through a sequence of stages. Each stage takes a defined input, applies a transformation, validates the result, and passes it forward. The important property of this architecture is determinism: the same input always produces the same output. That makes failures reproducible and gives each stage something concrete to depend on from the stage before it.
Most compiler pipelines separate planning from execution using an intermediate representation, a structured data format that sits between the two. The planner produces it, the compiler consumes it, and neither needs to understand the internals of the other. That separation keeps the pipeline modular and makes validation straightforward: constraints are checked at the intermediate stage, before any execution begins, so problems surface early.
In this system, WorldSpec is the intermediate representation. The language model takes a natural language prompt and produces a schema-locked WorldSpec. The compiler then transforms that WorldSpec into a runtime-executable world bundle. Walkability checks, asset budget enforcement, and spawn validation all happen at the WorldSpec stage, before any geometry loads into the headset. If something fails, the pipeline has a defined point to catch it and a deterministic path to recover from it.
\section{Related Work}
Research in prompt-driven 3D environment generation has split along one foundational question: should a system generate novel geometry from a prompt, or assemble environments from a curated library of pre-authored assets? That choice has consequences for reliability, performance, and what guarantees a downstream runtime can actually make.
Text2Room takes the generative geometry approach. Hollein et al. showed that a text prompt can produce a spatially consistent, textured 3D mesh by lifting outputs from 2D image generation models into a 3D representation~\cite{hoellein2023text2room}. The strength of this approach is coverage. Because geometry is synthesized rather than retrieved, the system can respond to arbitrary prompts without being limited by an asset catalog. In a live VR context, however, synthesized geometry requires post-generation analysis to determine walkability, draw call cost, and memory footprint before a runtime can safely load it. That additional step adds latency and introduces failure modes at the validation stage that are difficult to recover from within an interactive session. 3D-SceneDreamer addresses one dimension of this problem by representing global scene structure in a stable tri-plane NeRF and refining content in stages~\cite{Zhang_2024_CVPR_SceneDreamer}. The staged approach improves cross-view consistency and reduces cumulative geometric error, and that architectural pattern, stabilizing global structure first then refining locally, maps onto the greybox-first compilation strategy this system uses. The difference is that geometry synthesis still produces outputs without known collision or performance properties, whereas asset-based compilation carries those properties from the start.

The asset-based approaches accept a narrower generative range in exchange for stronger guarantees. Deitke et al. and Yang et al. both work within pre-authored asset libraries, using structured representations and constraint satisfaction to produce environments whose physical properties are known before the world loads~\cite{Deitke_2022_NeurIPS, Yang_2024_CVPR}. Deitke et al. focus on procedural layout generation at scale, producing large collections of varied indoor environments for embodied AI training~\cite{Deitke_2022_NeurIPS}. Yang et al. connect natural language to asset selection and spatial placement directly, using a language model to convert prompts into structured scene descriptions that get compiled into interactive environments~\cite{Yang_2024_CVPR}. Both systems treat generation as a pipeline with explicit validation stages, and both produce environments reliable enough to train agents in. Both systems assume the asset library contains what the prompt requests. Asset coverage is an implicit prerequisite, not a problem the pipeline resolves. Feng et al. extend this direction by showing that large language models can serve as visual planners, converting natural language prompts into stylesheet-like structured layouts through in-context demonstrations~\cite{Feng_2023_NeurIPS_LayoutGPT}. Their results support a clear separation of concerns: planning in structured symbolic space first, then compilation afterward. That separation directly motivates the WorldSpec stage in this system, where the language model produces a schema-locked intermediate representation rather than geometry, and a deterministic compiler handles the rest.

A related body of work approaches structured scene representation through scene graphs, where objects are nodes and spatial relationships are edges. GraphDreamer uses scene graph conditioning to generate compositional 3D scenes, introducing object-wise disentanglement and relationship-aware constraints to reduce interpenetration and improve spatial coherence~\cite{Gao_2024_CVPR_GraphDreamer}. CommonScenes takes a similar position, combining a layout branch and a shape generation branch conditioned on scene-object and object-object relationships, and reports improved semantic consistency compared with retrieval-heavy alternatives~\cite{Zhai_2023_NeurIPS_CommonScenes}. Both works provide evidence that explicit relational structure, represented before rendering rather than inferred from it, produces more controllable and reproducible outputs. For the purposes of this system, the scene graph functions as an analogue to WorldSpec: a machine-checkable intermediate representation that the compiler can validate against spatial and budget constraints before any geometry is loaded.

All three systems treat generation as an offline or training-time problem, where compilation time and frame stability fall outside the design requirements. A user standing in a headset waiting for a portal to open operates under a different set of conditions. The system described in this paper sits within the asset-based tradition of Deitke et al. and Yang et al., adopting structured intermediate representations and constraint-driven compilation for the same reliability reasons, but applies them to a live interactive setting where those guarantees have to hold in real time~\cite{Deitke_2022_NeurIPS, Yang_2024_CVPR}. A deterministic fallback planner handles cases where the language model produces invalid output. A greybox-first compilation phase ensures a walkable world exists before visual assets begin streaming. A substitution system replaces missing assets with the closest tagged alternatives, keeping the pipeline moving regardless of asset coverage gaps. The contribution is a composition of these properties into a pipeline that maintains reliability guarantees under the constraints of a standalone VR headset.
\section{Results}
Results will report world generation success rate, safe-spawn validity, prompt-to-world latency, and runtime stability across generated scenes.



\section{Conclusion}

\bibliographystyle{acm}
\bibliography{bibliography}

\end{document}
